\chapter{Conclusiones y Discusión Grupal}

A partir de la experiencia de laboratorio realizada, se presentan las siguientes conclusiones integradoras. El desarrollo de los ensayos permitió profundizar en la naturaleza de la medición física y la importancia del instrumental.

En primera instancia, la contrastación del voltímetro analógico puso de manifiesto fenómenos metrológicos críticos. La realización de pasadas ascendentes y descendentes permitió identificar la histéresis del instrumento, un error sistemático derivado de la fricción mecánica y las propiedades elásticas de los resortes internos. Si bien la realización de múltiples series de pasadas podría mejorar la repetibilidad al reducir errores aleatorios, se determinó que dos pasadas constituyen el compromiso mínimo necesario para detectar comportamientos bidireccionales sin enmascarar errores sistemáticos.

La validez de la curva de contrastación obtenida se estima en un periodo de 6 a 12 meses, siempre que el instrumento no
sufra golpes mecánicos o sobrecargas eléctricas que obliguen a un nuevo procedimiento. Un hallazgo fundamental de este
experimento fue que el voltímetro UNIVO se encuentra fuera de su clase original ($6 \percent$ experimental frente al
$2,5 \percent$ de catálogo), lo que resalta la necesidad de contar con patrones de referencia confiables en el
laboratorio, tales como los multímetros Pro'sKit MT-1210, UNI-T UT33A+ o cajas de décadas de precisión como la Time
Electronics 1051 (ver Anexo \ref{anexo:hoja_datos}).

En cuanto a la medición de baja resistencia, el empleo del método de cuatro terminales (Kelvin) demostró su superioridad técnica al eliminar los errores introducidos por la resistencia de contacto y de los cables de prueba. Para este ensayo, se utilizó una fuente de corriente basada en el regulador lineal LM317, el cual mantiene una corriente constante mediante una tensión de referencia estable (esquema detallado en el Anexo \ref{anexo:lm317}). Sin embargo, la experiencia evidenció que el método es tan robusto como el instrumental que lo implementa: debido a fallas técnicas en los rangos bajos del multímetro patrón, se debió emplear un rango de $10 \A$ para medir corrientes del orden de los $100 \mA$. Esto resultó en una incertidumbre relativa extremadamente elevada ($503 \percent$), invalidando la precisión del resultado final y dejando una lección práctica sobre la importancia de la correcta selección de escalas.

Sobre este último punto, el análisis de la propagación de incertidumbre permitió comparar el método de suma lineal frente al de raíz cuadrada de la suma de cuadrados (RSS). Se concluyó que el enfoque RSS proporciona una estimación más realista y probable según la normativa GUM, ya que contempla la baja probabilidad de que todas las fuentes de error alcancen sus valores máximos simultáneamente.

Finalmente, la medición de la resistencia de puesta a tierra mediante el uso de paños húmedos en el pasillo exterior resultó ser una adaptación efectiva ante la falta de terreno natural. El valor obtenido ($0,79 \ohm$) cumple satisfactoriamente con los niveles de seguridad, aunque se reconoce la influencia de la resistencia de contacto del montaje. En conjunto, las prácticas de laboratorio permitieron comprender que la medición profesional trasciende la simple lectura de un display, exigiendo un análisis crítico de la trazabilidad, el estado del instrumental y la correcta interpretación de las hojas de datos.

